\documentclass[12pt]{article}
\usepackage[a4paper, margin=1in]{geometry}
\usepackage{amsmath, amssymb, booktabs, setspace}
\setstretch{1.2}
\parskip=0.7em

\title{Tutorial  --The Heckscher--Ohlin Model}
\author{ECON 308 -- International Economic Relations}
\date{Fall 2025}

\begin{document}
\maketitle

\noindent
Two countries, \textbf{Home (Canada)} and \textbf{Foreign (Mexico)}, produce two goods — \textbf{Computers (C)} and \textbf{Textiles (T)} — using two factors of production: \textbf{Capital (K)} and \textbf{Labor (L)}.  
Both countries share identical technologies and preferences. The production functions for each good are:

\[
Q_C = (K_C)^{0.6}(L_C)^{0.4}, \quad Q_T = (K_T)^{0.3}(L_T)^{0.7}
\]

\noindent
The total factor endowments are:

\begin{center}
\begin{tabular}{lcc}
\toprule
\textbf{Country} & \textbf{Capital (K)} & \textbf{Labor (L)} \\
\midrule
Canada (Home) & 600 & 300 \\
Mexico (Foreign) & 300 & 600 \\
\bottomrule
\end{tabular}
\end{center}

\noindent
Computers are \textbf{capital-intensive}, while Textiles are \textbf{labor-intensive}.  
Both goods are produced under perfect competition and constant returns to scale.

%============================================================
\section*{1. Factor Abundance}

\begin{enumerate}
    \item Compute the \( K/L \) ratio for each country. Which country is capital-abundant and which is labor-abundant?
    \item Suppose both countries have the same technology. If Canada doubles both its capital and labor, does its factor abundance classification change?
\end{enumerate}
\newpage
%============================================================
\section*{2. Factor Intensity}

\begin{enumerate}
    \item Using:
    \[
    K_C = 300, \; L_C = 100; \quad K_T = 150, \; L_T = 350
    \]
    compute the \(K/L\) ratio in each industry.
    \item Derive the cost-minimizing capital-labor ratio for a Cobb–Douglas production function and explain how the wage–rental ratio (\(w/r\)) affects input choice.
\end{enumerate}

%============================================================
\section*{3. Autarky Equilibrium and Relative Prices}

\begin{enumerate}
    \item Suppose:
    \[
    (w/r)_{Canada} = 1.5, \quad (w/r)_{Mexico} = 3.0
    \]
    Compute the implied \(K/L\) ratios for Computers and Textiles in each country using the cost-minimization condition.
    \item Which good is relatively cheaper in each country under autarky? Explain.
    \item Qualitatively describe the relative supply curves for both countries and indicate the direction of comparative advantage.
\end{enumerate}

%============================================================
\section*{4. Trade Pattern and Terms of Trade}

\begin{enumerate}
    \item Predict each country’s exports and imports once trade opens.
    \item If autarky prices are:
    \[
    (P_C/P_T)_{Canada} = 1.2, \quad (P_C/P_T)_{Mexico} = 2.2
    \]
    and world price after trade is \( (P_C/P_T)^W = 1.8 \), determine how each country’s \textbf{terms of trade} change and the welfare implications.
\end{enumerate}
\newpage
%============================================================
\section*{5. Factor Price Equalization and Real Incomes}

\begin{enumerate}
    \item Use the Stolper–Samuelson theorem to determine how trade affects real wages and real rental rates in both countries.
    \item Suppose in Canada after trade the price of Computers rises by 20\% and Textiles falls by 10\%. Which factor gains more, and why?
    \item Can both factors gain from trade in the same country? Under what conditions could this occur?
\end{enumerate}

%============================================================
\section*{6. Relative Supply and World Equilibrium (Advanced)}

\begin{enumerate}
    \item Let the PPF of each country be approximated by:
    \[
    Q_T = A - B Q_C^2
    \]
    where \(A\) and \(B\) differ by country.  
    If \(A_{Canada}=200, B_{Canada}=0.005, A_{Mexico}=200, B_{Mexico}=0.02\), derive the relationship between \(P_C/P_T\) and \(Q_C\).
    \item Explain how the world equilibrium relative price is determined.
    \item Discuss how an increase in Canada’s capital stock affects world prices.
\end{enumerate}

\end{document}
